% ============================================================================
%  07_autodiagnostico.tex — Núcleo de la Actividad 3.
% ============================================================================
\section{Autodiagnóstico}
\label{sec:autodiag}

El módulo \texttt{diagnostics} es la \textbf{capa de instrumentación
avanzada} que se interpone entre las lecturas crudas de los sensores y el
lazo de control de clima. En cada \texttt{diag\_tick()} ejecuta cinco
acciones encadenadas: lee los sensores monitorizados, valida cada lectura por
rango, mantiene una media móvil de las medidas válidas, decide la fuente
activa de temperatura y de humedad por separado, y reacciona ante el fallo
con una severidad escalonada que puede llegar a inhibir el ascensor.

\paragraph{Detección de fallos por límites de rango.}
\label{sub:rango}

A cada sensor monitorizado se le asocia una \emph{banda de operación
plausible} en su unidad física (tabla~\ref{tab:bandas}). Una lectura fuera
de la banda ---o un \texttt{NaN} en el caso del DHT22, que la librería
devuelve ante fallo del protocolo--- se interpreta como
\textbf{cortocircuito o circuito abierto} de la conexión del sensor. Para
la temperatura y la humedad, los límites se han elegido más estrechos que
el rango total del chip (−40..80°C y 0..$100$\,\% según
\cite{aosong_dht22}): así un valor físicamente posible pero implausible
para el entorno de la planta se trata también como avería. Para evitar disparos por ruido, se exige que
\texttt{DIAG\_FAULT\_SAMPLES} (3) lecturas consecutivas sean malas antes de
declarar la avería; igualmente se exigen 3 buenas seguidas antes de dar el
sensor por recuperado (anti-rebote tipo Schmitt sobre la salud del sensor).

\begin{table}[H]
  \centering
  \small
  \begin{tabular}{p{0.18\textwidth} p{0.20\textwidth} p{0.18\textwidth} p{0.36\textwidth}}
    \toprule
    \textbf{Magnitud} & \textbf{Banda válida} & \textbf{Sensor} & \textbf{Justificación} \\
    \midrule
    Temperatura  & -10–50 °C & DHT22 (pri/res)
                 & Banda operativa industrial (más estrecha que el rango total
                   del sensor, −40..80°C). \\
    Humedad      & 15–100 \% & DHT22 (pri/res)
                 & < 15 \% implausible en aire ambiente. \\
    Iluminación  & 1–20000 lx     & LDR
                 & Interior industrial; > 20 000 lx son exteriores. \\
    Corriente    & ≤ 500 mA & Pot.\ acond.
                 & Umbral de \emph{sobreconsumo} (§\ref{sub:corriente}). \\
    \bottomrule
  \end{tabular}
  \caption{Bandas de operación plausible y su justificación.}
  \label{tab:bandas}
\end{table}

\paragraph{Promediado de medidas (media móvil).}
\label{sub:avg}

Para reducir el ruido de medida~\cite{corona2014instrumentacion} todas las
lecturas válidas pasan por una \textbf{media móvil} de
\texttt{DIAG\_AVG\_WINDOW}=5 muestras antes de alimentar al control y al
\emph{display}. El búfer de cada magnitud es independiente y descarta las
muestras marcadas como malas, de modo que un fallo transitorio no contamina
la media. Cuando una magnitud conmuta de sensor (ver §\ref{sub:redund}), su
búfer se \textbf{vacía} para no mezclar lecturas de dos físicas distintas.

\subsection{Redundancia independiente de temperatura y humedad}\label{sub:redund}

El enunciado pide \emph{duplicar el sensor de temperatura para que, si uno
falla, se active el sensor de reserva}. En esta entrega se ha llevado un paso
más allá: el DHT22 tiene un elemento de medida \emph{independiente} para
temperatura y otro para humedad, por lo que en la práctica puede degradarse
una magnitud sin la otra. El módulo lo refleja con cuatro monitores de salud
independientes (\texttt{monPriT}, \texttt{monBckT}, \texttt{monPriH},
\texttt{monBckH}) y \emph{dos} indicadores de fuente activa
(\texttt{usingBackupT} y \texttt{usingBackupH}). El resultado es que el
sistema puede estar usando el DHT22 \emph{de reserva} para la temperatura y
al mismo tiempo el \emph{principal} para la humedad, o viceversa. La
conmutación es automática y reversible: cuando un sensor en fallo se
recupera, su monitor confirma 3 lecturas buenas seguidas y la fuente vuelve
al principal.

El proceso de selección es: si el principal está sano se usa siempre el
principal; si no, se usa la reserva siempre que esté sana; si ambos fallan
se mantiene la última fuente activa. Cuando la fuente cambia, se vacía el
búfer de la media móvil para no mezclar lecturas de sensores distintos.

\paragraph{Diagnóstico de actuadores por sensor de corriente.}\label{sub:corriente}

La parte de actuadores se cubre con un \textbf{sensor de corriente} en
GPIO\,39 (V\textsubscript{N}). En un montaje real correspondería a la salida
ya acondicionada de una resistencia \emph{shunt} amplificada; en Wokwi se
modela con un potenciómetro, pues no es posible variar dinámicamente el
consumo real de un LED. El ADC se promedia con la misma ventana móvil que el
resto de magnitudes ---imprescindible para no disparar falsas alarmas con el
ruido del ADC del ESP32--- y se compara con
\texttt{DIAG\_CURRENT\_MAX\_MA}=500 mA: si la corriente
promediada supera ese umbral, se declara \textbf{sobreconsumo} y la
severidad sube a \textsf{ALARMA}.

\subsection{Reacción escalonada: OK → AVISO → ALARMA}\label{sub:severidad}

La decisión de severidad combina los siete monitores de salud anteriores con
una jerarquía sencilla pero realista (tabla~\ref{tab:severidad}). La idea es
distinguir un fallo con respaldo (donde el sistema puede seguir operando) de
un fallo crítico que exige detener el ascensor.

\begin{table}[H]
  \centering
  \small
  \begin{tabular}{p{0.18\textwidth} p{0.40\textwidth} p{0.36\textwidth}}
    \toprule
    \textbf{Severidad} & \textbf{Condición} & \textbf{Reacción} \\
    \midrule
    \textsf{OK}     & Todos los sensores sanos y operando con el principal.
                    & Operación normal. \\
    \textsf{AVISO}  & Algún DHT principal falló pero hay reserva sana
                      (sólo en T, sólo en H, o en ambas).
                    & Conmuta a reserva, lo indica en LCD y CSV, el sistema
                      sigue operando; \emph{sin} alarma sonora. \\
    \textsf{ALARMA} & Fallo crítico sin reserva: ambos DHT en T,
                      o ambos en H, o LDR averiado, o sobreconsumo.
                    & El ascensor pasa al estado \textsf{ALARM} (cabina
                      congelada), el buzzer pita cada
                      1 s y el LCD muestra el banner del fallo.
                      Al desaparecer el fallo, el sistema vuelve a
                      \textsf{IDLE} automáticamente. \\
    \bottomrule
  \end{tabular}
  \caption{Severidad escalonada del autodiagnóstico.}
  \label{tab:severidad}
\end{table}

Cuando la severidad es \textsf{ALARMA}, \texttt{diagnostics} llama a
\texttt{elevator\_setAlarm(true)} ---una pequeña API que se ha añadido al
módulo \texttt{elevator} para que el autodiagnóstico pueda forzar y liberar
el estado \textsf{ALARM} de la FSM existente. La cola de plantas se conserva
durante la alarma, de modo que al despejarse el fallo el ascensor reanuda
los destinos pendientes sin que el operador deba reintroducirlos.
